\section*{ВВЕДЕНИЕ}
\addcontentsline{toc}{section}{ВВЕДЕНИЕ}

Булевые функции играют ключевую роль в симметричной криптографии,
используясь как функции замены в блочных шифрах и компоненты генераторов
ключевого потока. Криптостойкость таких систем напрямую зависит от свойств
булевых функций: сбалансированности, нелинейности, дифференциального
поведения и алгебраической иммунности. Многие из этих требований
противоречат друг другу, что существенно усложняет конструирование функций
с оптимальным набором свойств.

\textbf{Актуальность темы} определяется тем, что центральным нелинейным
элементом современных блочных шифров является блок подстановки (S-блок),
стойкость которого полностью определяется свойствами составляющих его
булевых функций. Несмотря на то, что S-блоки действующих стандартов
удовлетворяют установленным криптографическим требованиям, ряд
теоретических вопросов остаётся открытым. Разработка инструмента, совмещающего анализ булевых свойств и направленный поиск подстановок с заданными характеристиками, создаёт
основу для воспроизводимых экспериментальных исследований в данной области.

\textbf{Целью работы} является разработка программного инструмента для
анализа криптографических свойств булевых функций и генерации S-блоков, удовлетворяющих заданным критериям стойкости
к дифференциальному и линейному криптоанализу.

Для достижения цели поставлены следующие \textbf{задачи}:
\begin{itemize}
    \item \textbf{изучить} криптографические характеристики булевых функций,
        используемые для оценки стойкости S-блоков;;
    \item \textbf{разработать} модуль вычисления булевых характеристик
        произвольного S-блока;
    \item \textbf{оценить} работоспособность инструмента на S-блоках
        известных шифров и оценить качество генерируемых подстановок.
\end{itemize}

\section{Современные блочные шифры и роль S-блоков}
\subsection{Общая структура блочных шифров}

\emph{Блочным шифром} называется симметричный алгоритм шифрования,
преобразующий блок открытого текста фиксированной длины $n$ бит в
блок шифртекста той же длины под управлением секретного ключа~$K$.
Шифрование задаётся биективной функцией
$E_K: \mathbb{F}_2^n \to \mathbb{F}_2^n$, обратимой при знании
ключа~\cite{alferov2002}. Типичные размеры блока современных шифров —
64, 128 или 256~бит.

Современные блочные шифры строятся как \emph{итерированные}: блок
$n$-битного открытого текста подвергается многократному применению
одной и той же \emph{раундовой функции} $R$, использующей различные
\emph{раундовые ключи} $K_1, K_2, \dots, K_r$, полученные из основного
ключа $K$ по алгоритму \emph{расширения ключа}:
\[
    E_K(x) = R_{K_r} \circ R_{K_{r-1}} \circ \cdots \circ R_{K_1}(x).
\]
Раундовая функция в большинстве конструкций реализует принцип
К.~Шеннона: совмещение операций \emph{нелинейной замены}
(substitution), создающей \emph{путаницу} (confusion), и
\emph{линейного перемешивания} (permutation), обеспечивающего
\emph{рассеяние} (diffusion).

\textbf{Подстановочно-перестановочная сеть (SPN).} В каждом раунде
последовательно применяются три преобразования:
\begin{enumerate}
    \item наложение раундового ключа (как правило, операцией $\oplus$);
    \item \emph{нелинейный слой} — параллельное применение S-блоков
        к фрагментам блока;
    \item \emph{линейный слой} — биективное линейное преобразование
        над $\mathbb{F}_2^n$, перемешивающее выходы S-блоков.
\end{enumerate}
По SPN-схеме построены AES, «Кузнечик», Serpent и SM4. Стойкость SPN
к линейному и дифференциальному криптоанализу определяется,
во-первых, свойствами S-блока, во-вторых, способностью линейного слоя
обеспечивать достаточное число активных S-блоков в любой
характеристике.

\textbf{Сеть Фейстеля.} Альтернативная архитектура, в которой блок
делится на две равные половины $(L, R)$, и раундовое преобразование
имеет вид
\[
    (L', R') = (R,\ L \oplus F_{K_i}(R)),
\]
где $F_{K_i}$ — раундовая функция, содержащая нелинейные S-блоки.
Особенностью данной архитектуры является то, что функция $F$ не обязана
быть биективной: обратимость всего раунда обеспечивается структурой
схемы. По схеме Фейстеля построены Twofish и Camellia; модифицированная
схема Фейстеля используется в SM4.

\textbf{Роль S-блока.} Во всех современных архитектурах S-блок является
\emph{единственным} источником нелинейности раундовой функции. Все
прочие компоненты — наложение ключа, линейный слой, перестановки —
являются линейными или аффинными над $\mathbb{F}_2$. Следовательно,
именно свойства S-блока определяют стойкость шифра к семейству атак,
эксплуатирующих наличие нетривиальной линейной или алгебраической
структуры (см. главы~2–4). Этим обусловлен интерес к построению
S-блоков, удовлетворяющих установленным криптографическим критериям, —
центральная задача настоящей работы.

\subsection{Обзор современных блочных шифров}
\indent Рассмотрим применение S-блоков в шести широко
применяемых блочных шифрах: \textit{AES}, ГОСТ~Р~34.12-2015 («Кузнечик»), \textit{Serpent},
\textit{Twofish}, \textit{SM4} и \textit{Camellia}.

\subsubsection{\textit{AES}}

В 1997 году Национальный институт стандартов и технологий США
инициировал открытый конкурс на замену устаревшего стандарта DES.
По итогам многолетней публичной экспертизы в 2001 году алгоритм Rijndael,
разработанный бельгийскими криптографами Йоаном Дайменом
и Винсентом Рейменом, был принят в качестве федерального
стандарта шифрования AES~\cite{fips197}.

AES является блочным шифром с размером блока 128~бит и поддержкой ключей
длиной 128, 192 и 256~бит. Структура алгоритма представляет собой
подстановочно-перестановочную сеть (\textit{SPN}). Каждый из 10, 12 или 14 раундов
(в зависимости от длины ключа) включает четыре преобразования:
\textit{SubBytes}, \textit{ShiftRows}, \textit{MixColumns} и \textit{AddRoundKey}.

AES использует единственный статичный 8-битный S-блок, представляемый
матрицей $16 \times 16$. Имеет нелинейность~112 при теоретическом максимуме~120, дифференциальную
равномерность~4 и алгебраическую степень~7.

% --- 2.2 Кузнечик ---
\subsubsection{ГОСТ~Р~34.12-2015}

Российский национальный стандарт блочного шифрования ГОСТ~Р~34.12-2015
был принят в 2015 году. Алгоритм «Кузнечик» разработан ФСБ России совместно
с компанией «ИнфоТеКС» и пришёл на смену ГОСТ~28147-89 («Магма»), чьи
S-блоки изначально имели засекреченное происхождение.~\cite{gost2015}.

«Кузнечик» оперирует блоками размером 128~бит при длине ключа 256~бит
и включает 10 раундов. Структура представляет собой подстановочно-перестановочную сеть.
Стандарт определяет один статичный 8-битный S-блок~$\pi$, заданный в виде
константной таблицы подстановки без раскрытия метода его построения.
Данный подход стал предметом научной дискуссии: в работе~\cite{biryukov2016}
было показано, что S-блок «Кузнечика» может быть представлен в виде
композиции определённых степеней в $\mathrm{GF}(2^8)$, что, по мнению ряда
исследователей, потенциально указывает на наличие скрытой структуры.
Нелинейность S-блока составляет~100, дифференциальная равномерность 8.
По обоим показателям он уступает S-блоку \textit{AES}, однако свойства линейного слоя
стандарта компенсируют данное различие на уровне всего шифра.

\subsubsection{Serpent}

\textit{Serpent} был представлен в 1998 году как один из пяти финалистов конкурса
AES. Авторами алгоритма являются Росс Андерсон (Великобритания), Эли Бихам
(Израиль) и Ларс Кнудсен (Дания). По итогам голосования \textit{Serpent} занял второе
место, набрав наибольшее число голосов среди участников, ставивших приоритетом
максимальную безопасность в ущерб производительности~\cite{serpent1998}.

Serpent использует блок размером 128~бит и ключи длиной 128, 192 или 256~бит.
Алгоритм включает 32 раунда, что вдвое превышает количество раундов \textit{AES}.
Авторы сознательно придерживались консервативного подхода, обеспечивающего
значительный запас стойкости к будущим атакам. Эффективность программных
реализаций достигается за счёт техники «битовых срезов» (\textit{bitslicing}),
позволяющей обрабатывать несколько блоков параллельно.

Serpent использует набор из восьми различных 4-битных S-блоков, применяемых
циклически по раундам. Каждый из них реализует биективное отображение
$\{0,1\}^4 \to \{0,1\}^4$. S-блоки выбирались из перестановок, производных
от S-блоков DES, с оптимизацией по критериям дифференциальной и линейной
стойкости. Нелинейность каждого S-блока равна~4~-максимально достижимой
для 4-битных булевых функций. Благодаря 32 раундам совокупная стойкость
\textit{Serpent} к дифференциальному и линейному криптоанализу превосходит
соответствующие показатели AES.

\subsubsection{Twofish}

\textit{Twofish} был разработан в 1998 году командой под руководством Брюса Шнайера
(Counterpane Internet Security) при участии Джона Келси, Дага Вина,
Дэвида Вагнера и Криса Холла. Алгоритм вошёл в пятёрку финалистов конкурса
\textit{AES}~\cite{twofish1998}.

\textit{Twofish} оперирует блоками размером 128~бит при ключах 128, 192 и 256~бит,
выполняя 16 раундов в структуре сети Фейстеля. Отличительной особенностью
алгоритма является использование ключезависимых S-блоков: конкретные таблицы
подстановки формируются из ключевого материала на этапе \textit{KeySchedule}.
Кроме того, \textit{Twofish} применяет псевдо-адамарово преобразование (PHT)
и \textit{MDS}-матрицу в функции Фейстеля.

Шифр использует четыре 8-битных S-блока, каждый из которых строится как
композиция двух фиксированных 4-битных перестановок ($q_0$ и $q_1$)
с промежуточным наложением ключевого материала через XOR и умножением на
\textit{MDS}-матрицу. Таким образом, алгебраическая структура S-блоков является
фиксированной, тогда как конкретные значения таблиц зависят от ключа.
Такой подход представляет собой компромисс между полностью статичными
и полностью случайными ключезависимыми подстановками и теоретически затрудняет
атаки, предполагающие предварительное знание S-блоков.

\subsubsection{SM4}

\textit{SM4} — китайский государственный стандарт блочного шифрования. Алгоритм
был разработан Государственным бюро криптографии КНР (\textit{OSCCA}) и впервые
опубликован в 2006 году в контексте стандарта беспроводных локальных
сетей \textit{WLAN}. В 2016 году \textit{SM4} был принят в качестве национального
стандарта и активно продвигается как альтернатива \textit{AES}
в государственных и коммерческих системах Китая~\cite{sm4_2006}.

\textit{SM4} использует блок размером 128~бит при длине ключа 128~бит и выполняет
32 раунда в обобщённой структуре сети Фейстеля. Каждый раунд применяет
функцию~$T$, состоящую из нелинейного преобразования~$\tau$ (побайтовое
применение S-блока) и линейного преобразования~$L$. Алгоритм расширения
ключа построен по принципу, аналогичному основному шифру.


\textit{SM4} использует один статичный 8-битный S-блок, заданный таблицей подстановки.
Происхождение S-блока долгое время оставалось официально нераскрытым,
что порождало вопросы, схожие с дискуссией вокруг «Кузнечика». Последующий
криптографический анализ показал, что S-блок \textit{SM4} построен по принципу,
аналогичному \textit{AES}: вычисление мультипликативной инверсии в $\mathrm{GF}(2^8)$
с использованием иного неприводимого полинома в сочетании с аффинным
преобразованием. Нелинейность S-блока составляет~112, дифференциальная
равномерность~4, что идентично параметрам \textit{AES}.


\subsubsection{Camellia}

\textit{Camellia }разработан совместно японскими компаниями \textit{NTT} и \textit{Mitsubishi Electric}
в 2000 году. Алгоритм прошёл европейский конкурс \textit{NESSIE}
и японский конкурс \textit{CRYPTREC}, получив статус рекомендованного шифра.
\textit{Camellia} включён в стандарты \textit{TLS}, \textit{IPsec} и широко применяется
в японской государственной и коммерческой инфраструктуре~\cite{camellia2001}.

\textit{Camellia }использует блок размером 128~бит при ключах 128, 192 или 256~бит.
Структура~ — сеть Фейстеля с 18 раундами (для ключа 128~бит) или 24 раундами
(192/256~бит). Характерной особенностью является наличие
$\mathrm{FL}/\mathrm{FL}^{-1}$-функций, вставляемых каждые 6 раундов:
они образуют дополнительный нелинейный слой, существенно усложняющий полный
дифференциальный и линейный криптоанализ. Алгоритм изначально проектировался
с учётом эффективности как программных, так и аппаратных реализаций.


\textit{Camellia} использует четыре 8-битных S-блока: один базовый ($s_1$) и три
производных ($s_2$, $s_3$, $s_4$), полученных циклическим сдвигом выходного
байта базового S-блока на 1~бит. Базовый S-блок построен аналогично \textit{AES}:
инверсия в $\mathrm{GF}(2^8)$ с иными параметрами поля и аффинным
преобразованием с другой матрицей. Нелинейность составляет~112,
дифференциальная равномерность~4. Использование четырёх родственных
S-блоков вместо одного повышает диффузию и усложняет алгебраический анализ,
сохраняя при этом компактность аппаратной реализации.

\subsection{Сравнительный анализ S-блоков современных стандартов}
Рассмотренные блочные шифры демонстрируют различные подходы к построению
нелинейного слоя. В \textit{AES}, \textit{SM4} и «Кузнечике» используется
единственный статичный 8-битный S-блок. В \textit{Camellia} применяются
четыре родственные 8-битные подстановки, полученные из базового S-блока.
Алгоритм \textit{Serpent} использует набор из восьми различных 4-битных
S-блоков, циклически применяемых по раундам. В \textit{Twofish} реализован
ключезависимый подход, при котором конкретные таблицы подстановки
формируются из ключевого материала.

С точки зрения количественных криптографических характеристик
наблюдаются устойчивые закономерности. Для 8-битных S-блоков
современных стандартов характерна нелинейность 112
(\textit{AES}, \textit{SM4}, \textit{Camellia}), что близко к
теоретическому максимуму 120. В «Кузнечике» нелинейность составляет 100.
Для 4-битных S-блоков \textit{Serpent} нелинейность равна 4,
что является максимально возможным значением для данной размерности.

Дифференциальная равномерность 8-битных подстановок в \textit{AES},
\textit{SM4} и \textit{Camellia} равна 4, что является оптимальным
значением для биективных отображений $\{0,1\}^8 \to \{0,1\}^8$.
В «Кузнечике» данный показатель равен 8. Для 4-битных S-блоков
\textit{Serpent} максимальная дифференциальная вероятность также
соответствует оптимальным значениям для выбранной размерности.

Алгебраическая степень 8-битных S-блоков рассматриваемых стандартов,
как правило, равна 7, что близко к максимально достижимому значению
для биективных отображений данной размерности.

Таким образом, анализ действующих стандартов позволяет выделить
практические ориентиры для 8-битных S-блоков современного уровня
криптографической стойкости: нелинейность не ниже 100,
дифференциальная равномерность не выше 8 и алгебраическая степень
не ниже 7. Полученные значения демонстрируют практический уровень
требований к S-блокам современных стандартов.

\subsection{Постановка задачи исследования и обоснование разработки инструмента}
Проектирование нелинейных компонентов блочных шифров S-блоков
представляет собой самостоятельную криптографическую задачу, актуальность
которой определяется несколькими факторами. 

Для биективных булевых функций размерности $8 \times 8$ точные границы достижимых
значений нелинейности при ограниченной дифференциальной равномерности
до сих пор не установлены, что делает экспериментальный поиск в данном
пространстве теоретически значимым. 

S-блоки действующих стандартов (\textit{AES}, \textit{SM4}, \textit{Camellia}) построены на основе инверсии в конечном
поле $\mathrm{GF}(2^8)$ и обладают компактным алгебраическим описанием,
что потенциально упрощает алгебраический анализ; поиск подстановок без
явной алгебраической структуры при сопоставимых метриках представляет
самостоятельный интерес.

Дополнительную сложность представляет наличие теоретических
противоречий между криптографическими характеристиками булевых функций: оптимизация одной характеристики нередко приводит к ухудшению другой. В результате проектирование S-блока
представляет собой задачу многокритериальной оптимизации, в которой
требуется найти допустимый компромисс между совокупностью количественных
показателей.

Несмотря на наличие устоявшихся стандартов, отсутствует общедоступный
универсальный инструмент, позволяющий:
\begin{itemize}
    \item вычислять полный набор булевых характеристик произвольного S-блока;
    \item анализировать их устойчивость к криптоанализу;
    \item осуществлять поиск S-блоков с заданными пороговыми значениями.
\end{itemize}

Отсутствие такого инструментария затрудняет воспроизводимость
исследований, сопоставимость результатов и экспериментальное
проектирование новых подстановок.

В связи с изложенным целью настоящей работы является разработка
программного инструмента, совмещающего функции анализатора булевых
свойств S-блоков и генератора подстановок, удовлетворяющих заданным
криптографическим критериям. Пороговые значения характеристик
определяются с учётом практических ориентиров, выявленных при анализе
действующих стандартов.
\subsection*{ВЫВОДЫ}
В первой главе рассмотрены современные блочные шифры и проанализированы
используемые в них S-блоки. Проведён сравнительный анализ количественных характеристик S-блоков,
выделены практические ориентиры,характерные для стандартов современного уровня стойкости.
\addcontentsline{toc}{subsection}{ВЫВОДЫ}


\newpage
\section{Стойкость к линейному криптоанализу}
\begin{definition}
S-блок (substitution box, блок подстановки) — это функция, выполняющая нелинейное преобразование в блочных шифрах.
S-блок размера $n\times m$ — это отображение:
\[S : \mathbb{F}_2^n \to \mathbb{F}_2^m\]
\end{definition}
\subsection{Линейный криптоанализ}

Суть линейного криптоанализа заключается в линейной аппроксимации единственного нелинейного компонента блочного шифра — S-блока. Для определения корреляции входных и выходных битов рассматривают следующее выражение:
\[
    \alpha_1x_1 \oplus \alpha_2x_2 \oplus \dots \oplus \alpha_nx_n \oplus \beta_1y_1 \oplus \beta_2y_2 \oplus \dots \oplus \beta_my_m = 0,
\]
где $\alpha = (\alpha_1, \alpha_2, \dots, \alpha_n) \in \mathbb{F}_2^n$ — маска входных битов; \\
\indent$\beta = (\beta_1, \beta_2, \dots, \beta_m) \in \mathbb{F}_2^m$ — маска выходных битов; \\
\indent$x = (x_1, \dots, x_n)$ — входной вектор; \\
\indent$y = (y_1, \dots, y_m) = S(x)$ — выходной вектор.

Используя скалярное произведение, это выражение можно записать в компактной форме:
\[
    (\alpha, x) \oplus \beta \cdot y = 0,
\]
где $(a,b) = a_1b_1 \oplus a_2b_2 \oplus \dots \oplus a_kb_k$ — скалярное произведение над $\mathbb{F}_2$.

\begin{definition}
Таблица линейных аппроксимаций (Linear Approximation Table, LAT) S-блока $S$ — это таблица размера $2^n \times 2^m$, элемент которой с индексами $[\alpha, \beta]$ определяется как:
\[
    \text{LAT}[\alpha, \beta] = |\{x \in \mathbb{F}_2^n \mid (\alpha, x) \oplus (\beta, S(x)) = 0\}| - 2^{n-1}
\]
\end{definition}

Значение $\text{LAT}[\alpha, \beta]$ показывает отклонение от равномерного распределения (смещение). Для идеального S-блока математическое ожидание $\text{LAT}[\alpha, \beta] = 0$.
\subsection{Критерий стойкости к линейному криптоанализу}

Эффективность линейной аппроксимации характеризуется её смещением. Чем больше абсолютное значение $|\text{LAT}[\alpha, \beta]|$, тем лучше аппроксимация для атакующего.

\begin{definition}
Линейность S-блока $S$ определяется как:
\[
    L(S) = \max_{\alpha, \beta \in \mathbb{F}_2^n \setminus \{0\}} |\text{LAT}[\alpha, \beta]|
\]
\end{definition}

\textbf{Критерий стойкости:} S-блок является стойким к линейному криптоанализу, если его линейность $L(S)$ минимальна.

\subsection{Аппроксимация как булева функция}

Рассмотрим линейную аппроксимацию как булеву функцию. Для фиксированных масок $\alpha$ и $\beta$ определим:
\[
    f_{\alpha,\beta}(x) =  (\alpha, x)  \oplus (\beta, S(x)) 
\]

Это булева функция $f_{\alpha,\beta} : \mathbb{F}_2^n \to \mathbb{F}_2$. 

Вес этой функции:
\[
    w(f_{\alpha,\beta}) = |\{x \in \mathbb{F}_2^n \mid f_{\alpha,\beta}(x) = 1\}|
\]

Связь с LAT:
\[
    \text{LAT}[\alpha, \beta] = 2^n - w(f_{\alpha, \beta}) - 2^{n-1} =2^{n-1} - w(f_{\alpha,\beta})
\]

\subsection{Нелинейность булевой функции}

\begin{definition}
Расстояние Хэмминга между функциями $f, g : \mathbb{F}_2^n \to \mathbb{F}_2$ определяется как:
\[
    d_H(f, g) = |\{x \in \mathbb{F}_2^n \mid f(x) \neq g(x)\}| = w(f \oplus g)
\]
\end{definition}

\begin{definition}
Нелинейность булевой функции $f : \mathbb{F}_2^n \to \mathbb{F}_2$ — это минимальное расстояние Хэмминга от $f$ до множества всех аффинных функций:
\[
    NL(f) = \min_{l \in \mathcal{A}_n} d(f, l) = 2^{n-1} - \frac{1}{2}\max_{a \in \mathbb{F}_2^n}|\widehat{f}(a)|,
\]
где $\mathcal{A}_n$ — множество всех аффинных функций от $n$ переменных; \\
\indent $\widehat{f}(\omega)$ — преобразование Уолша-Адамара функции $f$.
\end{definition}

Для S-блока $S : \mathbb{F}_2^n \to \mathbb{F}_2^m$ можно определить его покомпонентную нелинейность:
\[
    NL(S) = \min_{\beta \in \mathbb{F}_2^m \setminus \{0\}} NL(f_\beta),
\]
где $f_\beta(x) = (\beta, S(x))$ — булева функция, соответствующая линейной комбинации выходных координат.

\begin{theorem}
Связь между линейностью и нелинейностью S-блока:
\[
    L(S) = 2^{n-1} - NL(S)
\]
\end{theorem}

\begin{proof}
Рассмотрим булеву функцию $f_{\alpha, \beta}(x) = (\alpha, x) \oplus (\beta, S(x)) = (\alpha, x) \oplus f_\beta(x)$, где $f_\beta(x) = (\beta, S(x))$ — координатная функция S-блока.

\textbf{Шаг 1.} Если $f_{\alpha, \beta}(x) = 0$, то $f_\beta(x) = (\alpha, x)$. Если $f_{\alpha, \beta}(x) = 1$, то $f_\beta(x) = (\alpha, x) \oplus 1$.

\textbf{Шаг 2.} Расстояние Хэмминга от $f_\beta$ до аффинных функций $(\alpha, x)$ и $(\alpha, x) \oplus 1$:
\begin{align*}
    d_H(f_\beta, (\alpha, x)) &= wt(f_\beta \oplus (\alpha, x)) = wt(f_{\alpha, \beta}), \\
    d_H(f_\beta, (\alpha, x) \oplus 1) &= 2^n - d_H(f_\beta, (\alpha, x)) = 2^n - wt(f_{\alpha, \beta}).
\end{align*}

\textbf{Шаг 3.} Используя соотношение $w(f_{\alpha, \beta}) = 2^{n-1} - LAT[\alpha, \beta]$, получаем:
\begin{itemize}
    \item Если $LAT[\alpha, \beta] > 0$, то $wt(f_{\alpha, \beta}) < 2^{n-1}$, следовательно,
    \[
        \min\{d_H(f_\beta, (\alpha, x)), d_H(f_\beta, (\alpha, x) \oplus 1)\} = wt(f_{\alpha, \beta}) = 2^{n-1} - LAT[\alpha, \beta].
    \]
    
    \item Если $LAT[\alpha, \beta] < 0$, то $wt(f_{\alpha, \beta}) > 2^{n-1}$, следовательно,
    \[
        \min\{d_H(f_\beta, (\alpha, x)), d_H(f_\beta, (\alpha, x) \oplus 1)\} = 2^n - wt(f_{\alpha, \beta}) = 2^{n-1} + LAT[\alpha, \beta].
    \]
\end{itemize}

В обоих случаях минимальное расстояние равно $2^{n-1} - |LAT[\alpha, \beta]|$.

\textbf{Шаг 4.} Нелинейность функции $f_\beta$ — это минимальное расстояние до всех аффинных функций:
\[
    NL(f_\beta) = \min_{\substack{g \in \mathcal{A}_n}} d_H(f_\beta, g) = \min_{\alpha \in \mathbb{F}_2^n} \left(2^{n-1} - |LAT[\alpha, \beta]|\right) = 2^{n-1} - \max_{\alpha \neq 0} |LAT[\alpha, \beta]|,
\]
где $\mathcal{A}_n$ — множество всех аффинных функций от $n$ переменных.

\textbf{Шаг 5.} Нелинейность S-блока — это минимум по всем нетривиальным координатным функциям:
\begin{equation*}
\begin{aligned}
    NL(S) = \min_{\beta \neq 0} NL(f_\beta) = \min_{\beta \neq 0} \left(2^{n-1} - \max_{\alpha \neq 0} |LAT[\alpha, \beta]|\right) = \\
    2^{n-1} - \max_{\substack{\alpha \neq 0 \\ \beta \neq 0}} |LAT[\alpha, \beta]| = 2^{n-1} - L(S).
\end{aligned}
\end{equation*}
\end{proof}

\textbf{Следствие:} Максимизация нелинейности S-блока эквивалентна минимизации его линейности, что обеспечивает стойкость к линейному криптоанализу.


\subsection{Распределение нелинейностей координатных функций}

Для оценки качества S-блока рассматривают нелинейности всех нетривиальных линейных комбинаций выходов:
\[
    \{NL(f_\beta) \mid \beta \in \mathbb{F}_2^m \setminus \{0\}\}
\]

Помимо минимального значения $NL(S) = \min_{\beta \neq 0} NL(f_\beta)$, 
полезно вычислить:

\begin{itemize}
    \item \textbf{Среднюю нелинейность:}
    \[
        \overline{NL}(S) = \frac{1}{2^m - 1} \sum_{\beta \in \mathbb{F}_2^m \setminus \{0\}} NL(f_\beta)
    \]
    
    \item \textbf{Количество координатных функций с минимальной нелинейностью:}
    \[
        N_{\min} = |\{\beta \in \mathbb{F}_2^m \setminus \{0\} \mid NL(f_\beta) = NL(S)\}|
    \]
\end{itemize}

\textbf{Интерпретация:}
\begin{itemize}
    \item Если $N_{\min}$ мало, то у атакующего меньше выбора хороших аппроксимаций
    \item Высокая средняя нелинейность $\overline{NL}(S)$ свидетельствует о равномерной стойкости
    \item Малая дисперсия нелинейностей указывает на сбалансированность S-блока
\end{itemize}

\subsection*{ВЫВОДЫ}
\addcontentsline{toc}{subsection}{ВЫВОДЫ}

Во второй главе рассмотрена защита S-блока от линейного криптоанализа
Мацуи. Введены понятия таблицы линейных аппроксимаций (LAT) и линейности
$L(S)$ как количественной меры подверженности S-блока линейной
аппроксимации. Аппроксимация представлена как булева функция, что
позволило связать криптографическую характеристику S-блока со
свойством её компонентных булевых функций — \emph{нелинейностью}.
Доказана теорема о связи $L(S) = 2^{n-1} - NL(S)$, из которой следует,
что максимизация нелинейности $NL(S)$ эквивалентна минимизации
линейности $L(S)$ и тем самым обеспечивает стойкость к линейному
криптоанализу. Введены вспомогательные характеристики — средняя
нелинейность $\overline{NL}(S)$ и количество компонент с минимальной
нелинейностью $N_{\min}$, — позволяющие оценить распределение
стойкости по координатным функциям. В качестве рабочего ориентира
принято $NL(S) \geq 112$, соответствующее уровню S-блоков
действующих стандартов.

\newpage
\section{Стойкость к дифференциальному криптоанализу}

\subsection{Дифференциальный криптоанализ}

Дифференциальный криптоанализ предложен Э.~Бихамом и А.~Шамиром в 1990~году
для атаки на DES-подобные шифры~\cite{biham1991}. В отличие от линейного
криптоанализа, эксплуатирующего корреляции отдельных битов, в данной атаке
исследуется распространение \emph{разностей} пар входных значений через
раундовое преобразование.

\begin{definition}
Пусть $X', X'' \in \mathbb{F}_2^n$ — пара входных векторов. Их
\emph{разностью} называется значение
\[
    \Delta X = X' \oplus X''.
\]
Аналогично определяется выходная разность $\Delta Y = S(X') \oplus S(X'')$.
\end{definition}

Для идеального равномерно случайного отображения при фиксированной
$\Delta X \neq 0$ распределение $\Delta Y$ равномерно:
$P(\Delta Y \mid \Delta X) = 2^{-n}$. Атакующий ищет пары
$(\Delta X, \Delta Y)$, для которых эта вероятность существенно выше
случайной. Такие пары называют \emph{дифференциалами} S-блока, а
последовательность разностей по раундам шифра — \emph{дифференциальной
характеристикой}.

\medskip
\textbf{Независимость от ключа.} В шифрах с операцией наложения раундового
ключа $W = X \oplus K$ ключ \emph{не влияет} на входную разность S-блока:
\[
    \Delta W = W' \oplus W'' = (X' \oplus K) \oplus (X'' \oplus K) = X' \oplus X''.
\]
Следовательно, распределение разностей S-блока определяется только видом
отображения~$S$ и может быть проанализировано без знания ключа.

\medskip
\textbf{Связь с производной по направлению.} Выходная разность совпадает со
значением производной по направлению $\Delta X$:
\[
    S'_a(x) = S(x) \oplus S(x \oplus a),
\]
поэтому изучение разностного поведения S-блока сводится к изучению
распределений его производных по всем ненулевым направлениям $a$.

\subsection{Таблица разностного распределения}

\begin{definition}
\emph{Таблицей разностного распределения} (Difference Distribution Table, DDT)
S-блока $S: \mathbb{F}_2^n \to \mathbb{F}_2^n$ называется таблица размера
$2^n \times 2^n$, элемент которой с индексами $[a, b]$ определяется как
\[
    \mathrm{DDT}[a, b] = |\{x \in \mathbb{F}_2^n \mid S(x) \oplus S(x \oplus a) = b\}|.
\]
\end{definition}

Строка $a$ таблицы описывает распределение значений производной $S'_a$.
Элемент $\mathrm{DDT}[a, b]$ — количество пар $(x, x \oplus a)$, дающих
выходную разность $b$; вероятность дифференциала $(a, b)$ равна
$\mathrm{DDT}[a, b] / 2^n$.

\begin{theorem}[Свойства DDT биективного S-блока]
\label{thm:ddt-properties}
Пусть $S$ — биективное отображение $\mathbb{F}_2^n \to \mathbb{F}_2^n$. Тогда
для таблицы $\mathrm{DDT}$ выполнены следующие свойства:
\begin{enumerate}
    \item $\displaystyle\sum_{b \in \mathbb{F}_2^n} \mathrm{DDT}[a, b] = 2^n$ для любого $a$;
    \item $\displaystyle\sum_{a \in \mathbb{F}_2^n} \mathrm{DDT}[a, b] = 2^n$ для любого $b$;
    \item $\mathrm{DDT}[a, b]$ чётно для всех $a, b$;
    \item $\mathrm{DDT}[0, 0] = 2^n$ и $\mathrm{DDT}[0, b] = 0$ при $b \neq 0$.
\end{enumerate}
\end{theorem}

\begin{proof}
\textbf{(1)} При фиксированном $a$ каждому $x \in \mathbb{F}_2^n$
соответствует ровно одно значение $S(x) \oplus S(x \oplus a)$, поэтому
суммарное количество элементов в строке равно $|\mathbb{F}_2^n| = 2^n$.

\textbf{(2)} Из биективности $S$: при фиксированных $a, b$ уравнение
$S(x) \oplus S(x \oplus a) = b$ при варьировании $a$ по всем значениям и
фиксированном $b$ имеет в сумме столько решений $(x, a)$, сколько различных
пар $(y_1, y_2)$ с $y_1 \oplus y_2 = b$, то есть $2^n$.

\textbf{(3)} Если $x$ — решение уравнения $S(x) \oplus S(x \oplus a) = b$,
то $x \oplus a$ — также его решение. Следовательно, решения разбиваются
на пары, и их число чётно.

\textbf{(4)} Из определения: $S(x) \oplus S(x) = 0$ для всех $x$, что
даёт $\mathrm{DDT}[0, 0] = 2^n$. Так как сумма по строке $0$ равна $2^n$,
остальные элементы строки нулевые.
\end{proof}

\textbf{Следствие.} «Идеальная» таблица с $\mathrm{DDT}[a, b] = 1$ при
$a \neq 0$ невозможна — это противоречит свойству чётности (свойство~3
теоремы~\ref{thm:ddt-properties}).

\subsection{Дифференциальная равномерность}

\begin{definition}
\emph{Дифференциальной равномерностью} S-блока $S$ называется величина
\[
    \delta(S) = \max_{a \in \mathbb{F}_2^n \setminus \{0\},\ b \in \mathbb{F}_2^n} \mathrm{DDT}[a, b].
\]
\end{definition}

Значение $\delta(S)$ ограничивает максимальную вероятность одношагового
дифференциала: $\max_{(a,b),\,a\neq 0} P(\Delta Y = b \mid \Delta X = a) = \delta(S) / 2^n$.

\textbf{Критерий стойкости.} S-блок является стойким к дифференциальному
криптоанализу, если его дифференциальная равномерность $\delta(S)$
минимальна; вероятность любого одношагового дифференциала при этом
ограничена сверху величиной $\delta(S) / 2^n$.

\medskip
\textbf{Сложность атаки.} Объём требуемых пар выбранного открытого текста
оценивается как
\[
    N_D \approx c / p_D,
\]
где $p_D$ — вероятность многораундовой дифференциальной характеристики,
а $c$ — константа порядка единицы. При независимости активных S-блоков
характеристики
\[
    p_D = \prod_{i=1}^{\gamma} \frac{\mathrm{DDT}_i[a_i, b_i]}{2^n} \leq \left(\frac{\delta(S)}{2^n}\right)^{\gamma},
\]
где $\gamma$ — число \emph{активных} S-блоков (с ненулевой входной
разностью) в характеристике. Минимизация $\delta(S)$ в сочетании с
проектированием линейного слоя, обеспечивающим большое $\gamma$,
является базовым принципом стойкости SPN-шифров к
дифференциальному криптоанализу.

\subsection{Минимально достижимые значения дифференциальной равномерности}

\begin{theorem}
Для биективного S-блока $S: \mathbb{F}_2^n \to \mathbb{F}_2^n$ справедливо
$\delta(S) \geq 2$.
\end{theorem}

\begin{proof}
По свойству~3 теоремы~\ref{thm:ddt-properties} все элементы DDT чётные.
По свойству~1 сумма элементов строки $a \neq 0$ равна $2^n > 0$,
поэтому хотя бы один её элемент строго положителен, а значит,
не меньше~2.
\end{proof}

S-блоки, на которых достигается равенство $\delta(S) = 2$, называются
\emph{почти совершенно нелинейными} (Almost Perfect Nonlinear, APN);
концепция введена К.~Нюберг~\cite{nyberg1993apn}.
Для нечётных $n$ биективные APN-функции существуют (например, функция
$x^3$ над $\mathbb{F}_{2^n}$). Для чётных $n$ известны лишь небиективные
APN-функции; при $n = 8$ вопрос о существовании биективных APN остаётся
открытой проблемой~\cite{carlet2021book}. На практике для биективных
S-блоков размерности $8 \times 8$ минимальным достижимым значением
является $\delta(S) = 4$.

\medskip
\textbf{Эталонные значения.} Соответствующие S-блоки промышленных стандартов
имеют следующие значения дифференциальной равномерности:

\begin{center}
\begin{tabular}{lc}
\hline
S-блок & $\delta(S)$ \\
\hline
AES (Rijndael)     & 4 \\
Camellia           & 4 \\
SM4                & 4 \\
ГОСТ Р 34.12-2015 («Кузнечик») & 8 \\
\hline
\end{tabular}
\end{center}

S-блоки, построенные на основе мультипликативной инверсии в $\mathbb{F}_{2^8}$
(AES, Camellia, SM4), достигают теоретического минимума $\delta = 4$;
S-блок «Кузнечика», построенный композицией нелинейных перестановок,
имеет $\delta = 8$, что соответствует более слабой границе вероятности
дифференциала ($2 \cdot 2^{-8}$ против $2^{-6}$).

\subsection*{ВЫВОДЫ}
\addcontentsline{toc}{subsection}{ВЫВОДЫ}

В третьей главе рассмотрена защита S-блока от дифференциального
криптоанализа Бихама–Шамира. Введено понятие таблицы разностного
распределения и доказаны её основные свойства для биективных
отображений. В качестве количественной меры стойкости введена
дифференциальная равномерность $\delta(S)$; показано, что эта величина
ограничивает снизу вероятность одношагового дифференциала и
оценивает требуемое число пар выбранного открытого текста. Установлено,
что для биективных S-блоков размерности $8 \times 8$ минимальным
достижимым значением является $\delta = 4$, что согласуется с
характеристиками S-блоков большинства современных шифров.


\newpage
\section{Алгебраические свойства S-блоков}

Помимо линейного и дифференциального криптоанализа, к S-блоку
предъявляются требования, защищающие от двух атак, использующих его
\emph{алгебраическую структуру}: атаки высоких порядков (Lai,
1994~\cite{lai1994higher}; Knudsen, 1995~\cite{knudsen1995truncated})
и алгебраической атаки (Courtois–Meier, 2003~\cite{courtois2003algebraic}).
Соответствующие количественные критерии — \emph{алгебраическая степень}
и \emph{алгебраическая иммунность} компонентных булевых функций
S-блока — формулируются через алгебраическую нормальную форму.

\subsection{Алгебраическая нормальная форма и алгебраическая степень}

\begin{definition}
\emph{Алгебраической нормальной формой} (АНФ, полином Жегалкина) булевой
функции $f \in \mathcal{F}_n$ называется её представление в виде
XOR-суммы мономов:
\[
    f(x_1, \dots, x_n) = \bigoplus_{a \in \mathbb{F}_2^n} g(a) \cdot x_1^{a_1} x_2^{a_2} \cdots x_n^{a_n},
\]
где $g(a) \in \mathbb{F}_2$ — коэффициент при мономе, индексированном
вектором $a$; $x_i^0 = 1$, $x_i^1 = x_i$.
\end{definition}

АНФ существует и единственна для любой $f \in \mathcal{F}_n$; её
вычисление возможно через преобразование Мёбиуса за время
$O(n \cdot 2^n)$. Подробное изложение содержится в первой курсовой
работе.

\begin{definition}
\emph{Степенью монома} $x_1^{a_1} \cdots x_n^{a_n}$ называется число
$\mathrm{wt}(a) = \sum_{i=1}^n a_i$. \emph{Алгебраической степенью}
функции $f$ называется наибольшая степень монома, входящего в её АНФ
с ненулевым коэффициентом:
\[
    \deg f = \max\{\mathrm{wt}(a) \mid g(a) \neq 0\}.
\]
\end{definition}

\begin{definition}
\emph{Алгебраической степенью} S-блока $S: \mathbb{F}_2^n \to \mathbb{F}_2^n$
называется минимум степеней его нетривиальных компонентных функций:
\[
    \deg(S) = \min_{\beta \in \mathbb{F}_2^n \setminus \{0\}} \deg(f_\beta), \qquad f_\beta(x) = (\beta, S(x)).
\]
\end{definition}

Использование минимума (а не максимума) объясняется тем, что атакующему
достаточно одной компонентной функции низкой степени; критерий
стойкости требует, чтобы \emph{все} компоненты имели высокую степень.

\begin{theorem}
\label{thm:deg-bound}
Для биективного S-блока $S: \mathbb{F}_2^n \to \mathbb{F}_2^n$
справедливо $\deg(S) \leq n - 1$.
\end{theorem}

\begin{proof}
Пусть $f_\beta(x) = (\beta, S(x))$ — произвольная компонентная функция,
$\beta \neq 0$. Поскольку $S$ биективно, при пробегании $x$ по
$\mathbb{F}_2^n$ значения $S(x)$ также пробегают весь $\mathbb{F}_2^n$.
Тогда вес функции $f_\beta$ равен
\[
    w(f_\beta) = |\{y \in \mathbb{F}_2^n \mid (\beta, y) = 1\}| = 2^{n-1},
\]
т.~е. $f_\beta$ \emph{сбалансирована}. По формуле прямого преобразования
Мёбиуса коэффициент при мономе максимальной степени $x_1 x_2 \cdots x_n$
равен $w(f_\beta) \bmod 2 = 0$, поэтому моном степени $n$ не входит в
АНФ ни одной компоненты, и $\deg(S) \leq n - 1$.
\end{proof}

Для биективных S-блоков размерности $8 \times 8$ из
теоремы~\ref{thm:deg-bound} следует ограничение $\deg(S) \leq 7$;
эталонные S-блоки AES, SM4, Camellia и «Кузнечика» достигают этой
границы. Целевое значение в настоящей работе — $\deg(S) = 7$.

\subsection{Атака высоких порядков}

Атака высоких порядков (higher-order differential attack), предложенная
Х.~Лаем в 1994~году~\cite{lai1994higher} и развитая
Л.~Кнудсеном~\cite{knudsen1995truncated}, является обобщением
дифференциального криптоанализа на производные более высокого
порядка.

\begin{definition}
\emph{Производной порядка $k$} булевой функции $f$ по направлениям
$a_1, a_2, \dots, a_k \in \mathbb{F}_2^n$ называется функция
\[
    \Delta_{a_1, \dots, a_k} f(x) = \Delta_{a_k} \bigl(\Delta_{a_{k-1}} (\cdots \Delta_{a_1} f)\bigr)(x),
\]
где $\Delta_a f(x) = f(x \oplus a) \oplus f(x)$ — производная по
направлению.
\end{definition}

\begin{theorem}
\label{thm:hod}
Пусть $\deg f = d$. Тогда для любых $a_1, \dots, a_k \in \mathbb{F}_2^n$
\[
    \deg(\Delta_{a_1, \dots, a_k} f) \leq d - k.
\]
В частности, при $k > d$ производная порядка $k$ тождественно равна
нулю; при $k = d$ — является константой.
\end{theorem}

\begin{proof}
Пусть $f$ содержит моном степени $d$. Производная по направлению
от монома $x_{i_1} x_{i_2} \cdots x_{i_d}$ — это многочлен степени
не выше $d - 1$, поскольку в результате раскрытия
$f(x \oplus a) \oplus f(x)$ моном наивысшей степени сокращается.
Применяя производную $k$ раз, получаем понижение степени на $k$.
\end{proof}

\textbf{Атака.} Зная $\deg(S) = d$, атакующий выбирает $k = d + 1$
линейно независимых направлений $a_1, \dots, a_{d+1}$ и формирует
$2^{d+1}$ открытых текстов вида
\[
    \{x \oplus c_1 a_1 \oplus \cdots \oplus c_{d+1} a_{d+1} \mid (c_1, \dots, c_{d+1}) \in \mathbb{F}_2^{d+1}\}.
\]
По теореме~\ref{thm:hod} соответствующая производная порядка $d + 1$
тождественно равна нулю, то есть XOR-сумма всех $2^{d+1}$ шифртекстов
имеет известное фиксированное значение. Это даёт уравнение для битов
ключа последнего раунда и существенно ускоряет восстановление по
сравнению с полным перебором.

\textbf{Критерий стойкости.} Алгебраическая степень $\deg(S)$ должна
быть как можно выше — стойкость к атаке высоких порядков растёт
экспоненциально по $\deg(S)$. Для биективных S-блоков $8 \times 8$
целевое значение — $\deg(S) = 7$.

\subsection{Аннигиляторы и алгебраическая иммунность}

\begin{definition}
\emph{Аннигилятором} булевой функции $f \in \mathcal{F}_n$ называется
ненулевая функция $g \in \mathcal{F}_n$, удовлетворяющая равенству
\[
    f(x) \cdot g(x) = 0 \quad \forall x \in \mathbb{F}_2^n.
\]
\end{definition}

Эквивалентная характеризация: $\mathrm{supp}(g) \subseteq \overline{\mathrm{supp}(f)}$,
т.~е. носитель $g$ содержится в множестве нулей $f$.

\begin{definition}
\emph{Алгебраической иммунностью} функции $f$ называется минимальная
степень аннигилятора либо самой $f$, либо её отрицания $f \oplus 1$:
\[
    AI(f) = \min \{\deg g \mid g \neq 0,\ f \cdot g = 0 \text{ или } (f \oplus 1) \cdot g = 0\}.
\]
\end{definition}

Симметризация по $f$ и $f \oplus 1$ существенна: алгебраическая атака
эксплуатирует оба типа аннигиляторов.

\begin{definition}
\emph{Алгебраической иммунностью} S-блока $S$ называется минимум
$AI(f_\beta)$ по всем нетривиальным компонентным функциям:
\[
    AI(S) = \min_{\beta \in \mathbb{F}_2^n \setminus \{0\}} AI(f_\beta).
\]
\end{definition}

\begin{theorem}
\label{thm:ai-bound}
Для любой $f \in \mathcal{F}_n$ выполнено
\[
    AI(f) \leq \left\lceil n/2 \right\rceil.
\]
\end{theorem}

\begin{proof}[Идея доказательства]
Число мономов степени не выше $d$ от $n$ переменных равно
$M_d = \sum_{i=0}^{d} \binom{n}{i}$. При $d = \lceil n/2 \rceil$ имеем
$M_d > 2^{n-1}$, что превышает мощность носителя любой непостоянной
функции (равную либо $w(f)$, либо $2^n - w(f)$, не превосходящих
$2^{n-1}$ для сбалансированных $f$). Следовательно, система линейных
уравнений «моном степени $\leq d$ обращается в нуль на носителе $f$»
имеет ненулевое решение, и аннигилятор такой степени существует.
Полное доказательство приведено в~\cite{meier2004algebraic,carlet2021book}.
\end{proof}

Для $n = 8$ из теоремы~\ref{thm:ai-bound} следует $AI \leq 4$.
Заметим, что эта граница относится к булевой функции; для
\emph{S-блока} $S$ величина $AI(S) = \min_\beta AI(f_\beta)$
является минимумом по 255 компонентным функциям и может быть
строго меньше верхней границы. Прямое вычисление показывает, что
у S-блоков действующих стандартов (AES, Camellia, SM4, «Кузнечик»)
выполняется $AI(S) = 3$ (см. главу~6). Целевое значение в
настоящей работе — $AI(S) \geq 3$, согласованное с уровнем
эталонных S-блоков.

\subsection{Алгебраическая атака}

Алгебраическая атака (Courtois–Meier, 2003~\cite{courtois2003algebraic};
развитие — Meier–Pasalic–Carlet, 2004~\cite{meier2004algebraic})
эксплуатирует возможность описать шифр системой полиномиальных
уравнений низкой степени над $\mathbb{F}_2$. Если такую систему удаётся
составить, она решается линеаризацией или алгоритмом
XL/F$_4$ существенно быстрее полного перебора ключа.

\textbf{Роль аннигилятора.} Пусть компонентная функция $f_\beta$
S-блока удовлетворяет уравнению $f_\beta(x) = y_\beta$, где $y_\beta$ —
наблюдаемый бит шифртекста. Умножение уравнения на аннигилятор $g$
функции $f_\beta \oplus y_\beta$ даёт
\[
    g(x) \cdot (f_\beta(x) \oplus y_\beta) = 0,
\]
что эквивалентно уравнению $g(x) = 0$ степени $\deg g$. Чем меньше
$\deg g$, тем ниже степень результирующей системы. Следовательно,
\emph{минимальная} степень аннигилятора, т.~е. $AI(f_\beta)$, и
определяет стойкость S-блока к данной атаке.

\textbf{Критерий стойкости.} Алгебраическая иммунность $AI(S)$ должна
быть максимально близка к верхней границе $\lceil n/2 \rceil$.
В настоящей работе принимается порог $AI(S) \geq 3$ как практический
компромисс, согласованный с возможностями локального поиска.

\medskip
\textbf{Связь с алгебраической степенью.} Метрики $\deg(S)$ и $AI(S)$
\emph{независимы}: высокая алгебраическая степень не гарантирует
высокую алгебраическую иммунность. Известны функции с
$\deg f = n - 1$, обладающие линейным аннигилятором ($AI(f) = 1$);
обратные примеры также существуют. Этим объясняется необходимость
\emph{обоих} критериев при проектировании S-блока — они защищают от
разных атак.

\subsection*{ВЫВОДЫ}
\addcontentsline{toc}{subsection}{ВЫВОДЫ}

В четвёртой главе рассмотрены алгебраические свойства компонентных
булевых функций S-блока и их связь с двумя криптоаналитическими
атаками. Введены понятия алгебраической нормальной формы,
алгебраической степени и алгебраической иммунности. Показано, что для
биективных S-блоков размерности $8 \times 8$ достижимы значения
$\deg(S) = 7$ и $AI(S) = 4$, причём оба критерия независимы и
защищают от различных атак: высокая $\deg(S)$ препятствует атаке
высоких порядков (Lai, 1994), а высокая $AI(S)$ — алгебраической атаке
(Courtois–Meier, 2003). В качестве рабочих ориентиров приняты
$\deg(S) = 7$ и $AI(S) \geq 3$.


\newpage
\section{Методы генерации S-блоков}

Объединяя результаты глав~2–4, можно сформулировать задачу
проектирования S-блока как многокритериальную задачу оптимизации:
требуется построить биективное отображение $S: \mathbb{F}_2^8 \to \mathbb{F}_2^8$,
одновременно удовлетворяющее установленным пороговым значениям
четырёх метрик стойкости. Размер пространства поиска и отсутствие
аналитического градиента целевой функции делают неприменимыми
точные методы; в работе используются \emph{метаэвристические}
алгоритмы локального поиска — восхождение к вершине (hill climbing)
и имитация отжига (simulated annealing)~\cite{kirkpatrick1983sa}.

\subsection{Формализация задачи оптимизации}

\textbf{Пространство поиска.} Множеством допустимых состояний является
множество всех биективных отображений $\mathbb{F}_2^8 \to \mathbb{F}_2^8$,
изоморфное симметрической группе $\mathrm{Sym}_{256}$. Его мощность
\[
    |\mathrm{Sym}_{256}| = 256! \approx 8{,}6 \cdot 10^{506},
\]
что на много порядков превосходит число возможных состояний любой
вычислительной системы. Прямой перебор невозможен.

\textbf{Окрестность и мутация.} В качестве оператора локального
перехода используется \emph{транспозиция} двух элементов таблицы
подстановки:
\[
    \mathrm{swap}(S, i, j): S[i] \leftrightarrow S[j], \qquad i, j \in \{0, 1, \dots, 255\},\ i \neq j.
\]
Окрестность $N(S)$ состоит из всех состояний, получаемых из $S$
одной транспозицией; её мощность равна $\binom{256}{2} = 32640$.
Транспозиция автоматически сохраняет биективность, что устраняет
необходимость отдельной проверки допустимости после каждого шага.

\medskip
\textbf{Свойство инволютивности.} Повторное применение транспозиции
с теми же индексами возвращает таблицу в исходное состояние:
$\mathrm{swap}(\mathrm{swap}(S, i, j), i, j) = S$. Это свойство
позволяет реализовать «бесплатный» откат шага без хранения
снимков состояния — в работе оно используется в алгоритмах
калибровки и в hill climbing.

\textbf{Целевая функция.} Качество состояния измеряется скалярной
функцией штрафов $F: \mathrm{Sym}_{256} \to \mathbb{R}_{\geq 0}$,
описанной в подразделе~5.2. Задача оптимизации формулируется как
\[
    S^* = \arg\min_{S \in \mathrm{Sym}_{256}} F(S),
\]
причём целевое значение $F(S) = 0$ соответствует одновременному
выполнению всех четырёх пороговых критериев стойкости.

\subsection{Целевая функция штрафов}

В качестве целевой функции используется \emph{пороговая} fitness-функция,
накладывающая штраф только на нарушение установленных в главах~2–4
пороговых значений:
\begin{equation}
\label{eq:fitness}
\begin{aligned}
F(S) =\ &w_1 \cdot \max\bigl(0,\ 112 - N(S)\bigr) + w_2 \cdot \max\bigl(0,\ \delta(S) - 4\bigr) \\
       &+ w_3 \cdot \max\bigl(0,\ 7 - \deg(S)\bigr) + w_4 \cdot \max\bigl(0,\ \tau^* - AI(S)\bigr),
\end{aligned}
\end{equation}
где $w_1, w_2, w_3, w_4 > 0$ — настраиваемые весовые коэффициенты,
$\tau^* = 3$ — целевое значение алгебраической иммунности.

\textbf{Свойства функции~\eqref{eq:fitness}.}
\begin{enumerate}
    \item $F(S) \geq 0$ для любого $S$.
    \item $F(S) = 0$ тогда и только тогда, когда все четыре метрики
        удовлетворяют пороговым требованиям; это служит \emph{критерием
        остановки} алгоритма.
    \item Каждое слагаемое равно нулю, как только соответствующий
        порог достигнут, и далее не препятствует оптимизации
        остальных метрик. Эта \emph{разделимость} существенна, так как
        криптографические метрики, как правило, конкурируют:
        улучшение нелинейности может ухудшить алгебраическую степень
        и наоборот.
\end{enumerate}

\textbf{Альтернативные подходы.} В литературе распространена
функция Кларка–Джейкоба–Степни~\cite{clark2005sboxes} вида
\[
    \mathrm{CJS}(S) = \sum_{\beta \neq 0} \sum_{\omega \neq 0} \bigl|\,|\widehat{f_\beta}(\omega)| - X\,\bigr|^R,
\]
где параметры $X, R$ подбираются экспериментально. Эта функция
оптимизирует \emph{равномерность} спектра Уолша компонентных функций,
не делая нелинейность явной целью. В главе~7 показано, что для
биективных S-блоков размерности $8 \times 8$ данный подход приводит
к ландшафту с минимумом в стороне от целевого значения $N = 112$,
тогда как функция~\eqref{eq:fitness} имеет глобальный минимум
$F = 0$ именно в точке выполнения всех критериев.

\subsection{Восхождение к вершине}

Простейший алгоритм локального поиска. Из текущего состояния $S$
случайно выбирается сосед $S' \in N(S)$; переход осуществляется
тогда и только тогда, когда $F(S') < F(S)$. Ухудшающие шаги
\emph{никогда} не принимаются.

\begin{algorithm}{Hill climbing}\hfill\\
$S \leftarrow S_0$ \\
\textbf{повторять} пока есть улучшения или не превышен лимит итераций: \\
\quad $S' \leftarrow$ случайный сосед $S$ \\
\quad \textbf{если} $F(S') < F(S)$, \textbf{то} $S \leftarrow S'$ \\
\textbf{вернуть} $S$
\end{algorithm}

\textbf{Свойства.} Алгоритм прост в реализации, не имеет
гиперпараметров и быстро сходится. Главный недостаток — застревание
в \emph{любом} локальном минимуме целевой функции: если в
окрестности $N(S)$ нет улучшающих соседей, дальнейший прогресс
невозможен.

\textbf{Роль в работе.} Hill climbing используется в работе как
базовый ориентир (\emph{baseline}), позволяющий оценить, насколько
имитация отжига превосходит чисто жадную стратегию. Кроме того,
в подразделе~5.7 hill climbing используется как \emph{вторая фаза}
двухэтапного алгоритма, доводящая результат SA до целевого порога
по нелинейности.

\subsection{Имитация отжига}

\textbf{Физическая метафора.} Алгоритм имитации отжига
(simulated annealing, SA) предложен С.~Киркпатриком, К.~Геллаттом
и М.~Векки в 1983~году~\cite{kirkpatrick1983sa}; независимая
формулировка дана В.~Черни~\cite{cerny1985sa}. Название отражает
аналогию с физическим процессом отжига твёрдого тела: при медленном
охлаждении расплав успевает занять состояние с минимальной
энергией (правильную кристаллическую решётку); при быстром —
застывает в состоянии с локально, но не глобально минимальной
энергией (аморфное состояние). Идея SA — копировать «медленное
охлаждение» в дискретной оптимизационной задаче.

\textbf{Правило Метрополиса.} Ключевое отличие от hill climbing:
ухудшающие переходы принимаются с вероятностью, зависящей от
параметра $T > 0$ — \emph{температуры} алгоритма:
\[
    P\bigl(\text{принять переход } S \to S'\bigr) =
    \begin{cases}
        1, & \Delta F < 0, \\
        \exp\!\left(-\dfrac{\Delta F}{T}\right), & \Delta F \geq 0,
    \end{cases}
\]
где $\Delta F = F(S') - F(S)$. При $T \to \infty$ принимается любой
переход (случайное блуждание); при $T \to 0$ алгоритм вырождается
в hill climbing.

\textbf{Расписание температуры.} Температура убывает по итерациям
согласно \emph{расписанию} $T(t)$. В работе используется
\emph{геометрическое} расписание
\[
    T_{n+1} = \alpha \cdot T_n, \qquad \alpha \in (0; 1),
\]
с типичными значениями $\alpha \in [0{,}95; 0{,}99]$. Логарифмическое
расписание $T_n = T_0 / \log(n+1)$, обеспечивающее теоретическую
сходимость к глобальному оптимуму~\cite{gemangeman1984sa},
на практике требует неприемлемо большого числа итераций.

\textbf{Внутренний цикл.} На каждой температуре выполняется $M$
шагов мутации до перехода к следующей температуре. Это позволяет
системе достигнуть приближённого \emph{теплового равновесия}
при заданной $T$ до её снижения.

\textbf{Best-seen.} Параллельно с текущим состоянием $S$ алгоритм
поддерживает лучшее найденное за всё время работы состояние
$S^*$: возвращается именно оно, так как из-за случайных
принятий ухудшений финальное состояние может оказаться хуже
ранее посещённых.

\begin{algorithm}{Simulated annealing}\hfill\\
\textbf{Вход:} $S_0$, $T_0$, $\alpha$, $M$, $T_{\min}$, $F_{\mathrm{target}}$ \\
\textbf{Выход:} $S^*$ --- лучшее найденное состояние \\[2pt]
$S \leftarrow S_0$, $S^* \leftarrow S$, $T \leftarrow T_0$ \\
\textbf{пока} $T > T_{\min}$: \\
\quad \textbf{для} $k = 1, \dots, M$: \\
\quad\quad $S' \leftarrow$ случайный сосед $S$ \\
\quad\quad $\Delta F \leftarrow F(S') - F(S)$ \\
\quad\quad \textbf{если} $\Delta F < 0$ или $\mathrm{rand}() < \exp(-\Delta F / T)$: \\
\quad\quad\quad $S \leftarrow S'$ \\
\quad\quad\quad \textbf{если} $F(S) < F(S^*)$: $S^* \leftarrow S$ \\
\quad\quad\quad \textbf{если} $F(S^*) \leq F_{\mathrm{target}}$: \textbf{вернуть} $S^*$ \\
\quad $T \leftarrow \alpha \cdot T$ \\
\textbf{вернуть} $S^*$
\end{algorithm}

\textbf{Общее число итераций.} При геометрическом расписании
\[
    N_{\mathrm{total}} = M \cdot \frac{\ln(T_{\min}/T_0)}{\ln \alpha}.
\]
Например, при $T_0 = 100$, $T_{\min} = 10^{-3}$, $\alpha = 0{,}99$,
$M = 10^3$ получаем $N_{\mathrm{total}} \approx 1{,}15 \cdot 10^6$ шагов.

\subsection{Калибровка начальной температуры}

Начальная температура $T_0$ — один из наиболее критичных
гиперпараметров SA, поскольку \emph{зависит от масштаба} целевой
функции, заранее неизвестного. Слишком малая $T_0$ превращает SA
в hill climbing; слишком большая — в случайное блуждание на
протяжении многих итераций.

\textbf{Идея калибровки.} Вместо непосредственного выбора $T_0$
задаётся целевая доля $p_0$ ухудшающих переходов, которые должны
приниматься в начале работы (типично $p_0 = 0{,}8$). Из правила
Метрополиса
\[
    p_0 = \exp\!\left(-\frac{\langle\Delta F\rangle}{T_0}\right)
\]
выражается
\begin{equation}
\label{eq:t0}
    T_0 = -\frac{\langle\Delta F\rangle}{\ln p_0},
\end{equation}
где $\langle\Delta F\rangle$ — среднее значение $\Delta F$ по
ухудшающим переходам, оценённое эмпирически (Ben-Ameur, 2004~\cite{benameur2004t0}).

\begin{algorithm}{Оценка $T_0$}\hfill\\
\textbf{Вход:} $S_0$, $K$ — число пробных шагов, $p_0$ — целевая acceptance rate \\
\textbf{Выход:} $T_0$ \\[2pt]
$\Sigma \leftarrow 0$, $c \leftarrow 0$ \\
\textbf{для} $k = 1, \dots, K$: \\
\quad $(i, j) \leftarrow$ случайные индексы; $\Delta F \leftarrow F(\mathrm{swap}(S_0, i, j)) - F(S_0)$ \\
\quad \textbf{если} $\Delta F > 0$: $\Sigma \leftarrow \Sigma + \Delta F$, $c \leftarrow c + 1$ \\
\quad $\mathrm{swap}(S_0, i, j)$ \quad // откат \\
\textbf{если} $c = 0$: \textbf{вернуть} $1$ \\
\textbf{вернуть} $-\Sigma / (c \cdot \ln p_0)$
\end{algorithm}

\textbf{Откат каждого пробного шага} существенен: без него
состояние $S_0$ дрейфует во время калибровки, и собранная
статистика перестаёт характеризовать стартовую точку. Свойство
инволютивности транспозиции обеспечивает откат за константное
время.

\subsection{Multi-start и параллельные запуски}

Имитация отжига не гарантирует нахождение глобального минимума
за конечное число итераций: один запуск алгоритма представляет
собой одну выборку из распределения возможных финальных
состояний. Стратегия \emph{multi-start} состоит в выполнении $N$
независимых запусков с различными случайными начальными
состояниями и возвращении лучшего из найденных.

\textbf{Вероятностная оценка.} Если вероятность одного запуска
найти решение качества не хуже заданного равна $p$, то
вероятность хотя бы одного успешного запуска среди $N$
независимых:
\begin{equation}
\label{eq:multistart}
    P_N = 1 - (1 - p)^N.
\end{equation}
Например, при $p = 0{,}1$ десять запусков дают $P_{10} \approx 0{,}65$,
двадцать запусков — $P_{20} \approx 0{,}88$, пятьдесят — более
$0{,}99$. Для редко достижимых целей multi-start является
необходимой надстройкой над одиночным SA~\cite{marti2013multistart}.

\textbf{Параллелизация.} Запуски не имеют общего состояния и не
обмениваются информацией, что относит multi-start к классу
\emph{embarrassingly parallel} задач. На вычислителе с $K$
процессорными ядрами $N$ запусков выполняются за время, близкое
к $\lceil N/K \rceil$ времени одного запуска; при $N = K$ полное
расписание укладывается во время одного запуска. В реализации
используется $N = 16$ параллельных запусков, что соответствует
числу логических ядер используемого вычислителя.

\textbf{Калибровка $T_0$ в multi-start.} Поскольку $T_0$ зависит
от стартовой точки, в общем случае калибровку
по~\eqref{eq:t0} необходимо выполнять отдельно для каждого
запуска. На практике распределение средних $\langle\Delta F\rangle$
по случайным биективным S-блокам имеет узкий разброс (см. главу~7),
поэтому допустимо использовать общее значение $T_0$, оценённое
по одной стартовой точке.

\subsection{Двухфазный алгоритм SA + HC по нелинейности}

Экспериментальные данные (см. главу~7), а также результаты
литературы~\cite{clark2005sboxes,mdpi2022sboxes} показывают, что
\emph{одиночная} имитация отжига для биективных S-блоков
размерности $8 \times 8$ стабильно достигает значения нелинейности
$N \in [100; 104]$ и далее не прогрессирует. Этот барьер не
устраняется ни увеличением числа итераций, ни модификацией
cost-функции и носит структурный характер.

\textbf{Двухфазная схема.} Кларк, Джейкоб и Степни в работе~\cite{clark2005sboxes}
предложили следующую схему, известную как
SNLT (Simulated Nonlinearity Targeting):
\begin{enumerate}
    \item \textbf{Фаза 1 — глобальный поиск} (SA с пороговой
        fitness-функцией): начиная со случайной биективной
        перестановки, найти S-блок с нелинейностью $N \approx 100$–$104$.
        Цель фазы — выход из «начальной» области ландшафта в окрестность
        бассейна притяжения целевого значения.
    \item \textbf{Фаза 2 — локальное доуточнение} (HC по $N$):
        исходя из выхода фазы~1, выполнить hill climbing
        \emph{непосредственно по нелинейности}, принимая
        транспозицию тогда и только тогда, когда она не уменьшает
        $N$. Цель — доведение значения с $\approx 104$ до целевого
        $112$.
\end{enumerate}

\begin{algorithm}{Двухфазный SA + HC}\hfill\\
\textbf{Вход:} расписание $\sigma$, fitness $F$, число шагов HC --- $H$ \\
\textbf{Выход:} $S^*$ \\[2pt]
$S \leftarrow \mathrm{SA}(S_0, \sigma, F)$ \quad // фаза 1 \\
\textbf{для} $h = 1, \dots, H$: \\
\quad $(i, j) \leftarrow$ случайные индексы \\
\quad $\mathrm{swap}(S, i, j)$ \\
\quad \textbf{если} $N(S) < N_{\mathrm{prev}}$: $\mathrm{swap}(S, i, j)$ \quad // откат \\
\quad \textbf{иначе}: $N_{\mathrm{prev}} \leftarrow N(S)$ \\
\textbf{вернуть} $S$
\end{algorithm}

\textbf{Интерпретация.} Hill climbing из случайной стартовой точки
достигает $N \approx 98$ и застревает; hill climbing из стартовой
точки SA, лежащей вблизи $N = 104$, оказывается ближе к
\emph{бассейну притяжения} целевого значения $N = 112$, и среди
соседей с большей вероятностью находятся улучшающие переходы.
Эмпирическим обоснованием служит таблица~III статьи Кларка–Джейкоба–Степни,
в которой двухфазный алгоритм достигает $N = 112$, в то время как
одиночные SA и HC устойчиво останавливаются в районе
$N = 100$–$104$.

В сочетании с multi-start подходом из подраздела~5.6 двухфазный
алгоритм запускается параллельно $N$ раз; финальное решение
выбирается как наилучшее по фитнесу из всех ветвей.

\subsection*{ВЫВОДЫ}
\addcontentsline{toc}{subsection}{ВЫВОДЫ}

В пятой главе сформулирована задача проектирования S-блока как
задача минимизации скалярной пороговой fitness-функции на
пространстве биективных подстановок и описаны методы её решения.
Подробно изложены алгоритмы восхождения к вершине и имитации отжига;
доказана связь между ними через предельный переход $T \to 0$.
Описаны два технических приёма, существенных для практической
эффективности SA: калибровка начальной температуры по целевой
acceptance rate и стратегия multi-start с параллельным выполнением
независимых запусков. Сформулирован двухфазный алгоритм SA + HC,
позволяющий преодолеть структурный барьер одиночного SA и
достичь целевого значения нелинейности $N = 112$. Конкретные
параметры алгоритмов и эмпирические результаты их применения
приводятся в главе~7.


\newpage
\section{Программная реализация инструмента}

В рамках работы разработан программный инструмент анализа и
генерации S-блоков на языке Go. Выбор языка обусловлен тремя
соображениями: статической типизацией, повышающей надёжность
численного кода; компиляцией в собственный машинный код, дающей
производительность, сопоставимую с C; и встроенной моделью лёгких
потоков (\emph{горутин}), естественно отображающей задачу
параллельного multi-start (см. подраздел~5.6) на 16-ядерный
вычислитель, использованный для экспериментов.

\subsection{Архитектура и структура модулей}

Код разделён на четыре пакета по принципу возрастающей абстракции:

\begin{itemize}
    \item \texttt{boolfunc} — булевы функции от $n$ переменных:
        формы представления, преобразования, метрики;
    \item \texttt{sbox} — S-блок 8$\times$8 как биективная
        подстановка $\mathbb{F}_2^8 \to \mathbb{F}_2^8$:
        компонентные функции, DDT, агрегированные метрики;
    \item \texttt{generate} — алгоритмы генерации (SA, hill climbing,
        multi-start) и целевые функции;
    \item \texttt{verify} — модуль верификации построенных S-блоков
        (на текущей стадии работы — заготовка).
\end{itemize}

Зависимости направлены от менее абстрактных пакетов к более:
\texttt{sbox} использует \texttt{boolfunc} для работы с компонентными
функциями; \texttt{generate} использует \texttt{sbox} для оценки
метрик в cost-функциях. Циклических зависимостей нет.

Эталонные S-блоки промышленных шифров (AES, Camellia, «Кузнечик»,
SM4) хранятся в каталоге \texttt{testdata/} в формате JSON и
используются как тестовые векторы для модульного тестирования и
как контрольные точки для верификации правильности реализованных
метрик.

\subsection{Модуль булевых функций}

Тип \texttt{BoolFunc} представляет булеву функцию $f: \mathbb{F}_2^n \to \mathbb{F}_2$
её \emph{вектором истинности} (TV) — массивом длины $2^n$, где
$\mathrm{tv}[i] = f(i)$. Все вычисляемые характеристики
(вес, АНФ, спектр Уолша, нелинейность, степень, иммунность)
\emph{ленивы} и кэшируются на уровне поля структуры: повторный
вызов соответствующего геттера возвращает сохранённое значение
без перевычисления.

Дополнительно поддерживается \emph{упакованное} представление
TV в виде массива слов \texttt{[]uint64}: 64 значения функции в
одном машинном слове. Это даёт оценку веса через инструкцию
\texttt{POPCNT} над всем словом за один такт.

\textbf{Конструкторы.} Реализованы четыре способа создания
\texttt{BoolFunc}:
\begin{itemize}
    \item \texttt{New(s string)} — из строки нулей и единиц;
    \item \texttt{FromSlice(tv []uint8)} — из готового вектора;
    \item \texttt{Random(n)} — случайная функция от $n$ переменных;
    \item \texttt{RandomBalanced(n)} — случайная \emph{сбалансированная}
        функция (ровно $2^{n-1}$ единиц).
\end{itemize}
Все конструкторы валидируют входные данные: длина TV должна быть
положительной степенью двойки, элементы — лежать в $\{0, 1\}$.

\textbf{Преобразования.} Два центральных преобразования реализованы
in-place в общем виде с асимптотикой $O(n \cdot 2^n)$:

\emph{Преобразование Мёбиуса} (вычисление АНФ) реализовано
«бабочкой» по $n$ слоям: на $i$-м слое для каждого блока размером
$2^{i+1}$ выполняется XOR верхней половины блока с нижней. Это
прямой аналог быстрого преобразования Уолша–Адамара (FWHT) над
$\mathrm{GF}(2)$.

\emph{Быстрое преобразование Уолша–Адамара} реализовано на знаковом
векторе $\mathrm{sv}[i] = (-1)^{f(i)} \in \{-1, +1\}$ с
коэффициентами $W_f(a) = \sum_x (-1)^{f(x) \oplus \langle a, x \rangle}$.
Возвращаемый массив имеет тип \texttt{[]int32} (значения в $[-2^n; 2^n]$
помещаются в 32 бита с большим запасом).

\textbf{Метрики.} В таблице~\ref{tab:boolfunc-metrics} приведены
основные методы модуля и сложность их первого вызова.

\begin{table}[h]
\centering
\caption{Метрики \texttt{BoolFunc}: формулы и сложность.}
\label{tab:boolfunc-metrics}
\begin{tabular}{lll}
\hline
Метод & Формула & Сложность \\
\hline
\texttt{Weight()}           & $|\{x : f(x) = 1\}|$            & $O(2^n / 64)$ \\
\texttt{ANF()}              & $\mu(\mathrm{tv})$              & $O(n \cdot 2^n)$ \\
\texttt{Walsh()}            & $\widehat{f}$ через FWHT        & $O(n \cdot 2^n)$ \\
\texttt{Nonlinearity()}     & $2^{n-1} - \tfrac{1}{2}\max|\widehat{f}|$ & $O(n \cdot 2^n)$ \\
\texttt{AlgebraicDegree()}  & $\max\{\mathrm{popcount}(i) : g(i) = 1\}$ & $O(2^n)$ + ANF \\
\texttt{AlgebraicImmunity()}& алгоритм гауссова исключения над $\mathbb{F}_2$ & $O(M_d \cdot 2^n)$ \\
\hline
\end{tabular}
\end{table}

\textbf{Вычисление алгебраической иммунности.} Алгоритм проверки
существования аннигилятора степени не выше $d$ реализован
через построение матрицы значений всех мономов степени не выше
$d$ на носителе функции и вычисление её ранга над $\mathbb{F}_2$
методом гауссова исключения. Ранг, меньший числа столбцов,
свидетельствует о наличии ненулевого ядра, т.~е. ненулевого
аннигилятора (см. подраздел~4.3). Строки матрицы упакованы в
\texttt{uint64}, что позволяет одной операцией XOR над словом
выполнять одновременное исключение 64~столбцов. Для определения
$AI(f)$ алгоритм выполняется последовательно для $d = 1, 2, \dots,
\lceil n/2 \rceil$ и для пары $(f, f \oplus 1)$ согласно
определению; первое значение $d$, на котором найден аннигилятор,
возвращается как результат.

\subsection{Модуль S-блока}

Тип \texttt{SBox} представляет S-блок $8 \times 8$ как массив
\texttt{[256]uint8}. Аналогично \texttt{BoolFunc} применяется ленивое
кэширование: компонентные функции, таблица DDT и агрегированные
метрики вычисляются при первом обращении и сохраняются в полях
структуры.

\textbf{Базовые проверки.} Метод \texttt{IsBijective()} проверяет,
что каждое значение из $\{0, \dots, 255\}$ встречается в таблице
ровно один раз. Методы \texttt{FixedPoints()} и
\texttt{OppositeFixedPoints()} возвращают, соответственно, множества
неподвижных ($S(x) = x$) и антиподвижных ($S(x) = \overline{x}$)
точек — оба критерия учитываются при отборе S-блоков в современных
стандартах.

\textbf{Компонентные функции.} Метод \texttt{Component(b)}
возвращает компонентную функцию $f_\beta(x) = \langle \beta, S(x) \rangle$
для маски $\beta = b \in \{1, \dots, 255\}$. Скалярное произведение
по модулю~2 вычисляется через свёртку: применяется маска
\texttt{b \& S[x]}, после чего значение сворачивается до одного бита
тремя XOR-сдвигами без ветвления. Результирующая
\texttt{BoolFunc} кэшируется и используется при вычислении
агрегированных метрик \texttt{Nonlinearity()},
\texttt{AlgebraicDegree()}, \texttt{AlgebraicImmunity()}.

\textbf{Таблица DDT.} Метод \texttt{DDT()} строит полную таблицу
$\mathrm{DDT}[a][b] = |\{x : S(x \oplus a) \oplus S(x) = b\}|$ за
один проход по $a \in \{1, \dots, 255\}$, $x \in \{0, \dots, 255\}$;
сложность $O(2^{2n}) = O(65536)$. Параллельно фиксируется максимум —
дифференциальная равномерность $\delta(S)$, возвращаемая методом
\texttt{DiffUniformity()}.

\textbf{Агрегированные метрики.} Поскольку нелинейность и
алгебраическая степень S-блока определяются как минимум/максимум
по 255 компонентным функциям, эти величины вычисляются за один
общий проход в методе \texttt{computeMetrics()}, существенно
амортизируя стоимость построения компонент.

\textbf{Загрузка из JSON.} Функция \texttt{LoadFromFile(path)}
читает эталонные S-блоки из каталога \texttt{testdata/}, где
подстановка хранится в виде матрицы $16 \times 16$
(стандартный формат представления S-блоков в технической
литературе). Это позволяет использовать общие тестовые векторы
для всех модулей и для верификации (см. подраздел~6.5).

\subsection{Модуль генерации}

\textbf{Состояние SA.} Тип \texttt{SAState} хранит текущую
подстановку, её закэшированное значение cost-функции и саму
функцию. Метод \texttt{Swap(i, j)} обменивает \texttt{s[i]} и
\texttt{s[j]} и пересчитывает cost; повторный вызов с теми же
индексами возвращает состояние в исходное — это и есть свойство
инволютивности, рассмотренное в подразделе~5.1. Метод
\texttt{RandomSwap()} выбирает пару различных индексов случайно и
возвращает старую cost для последующего вычисления $\Delta F$.

\textbf{Cost-функции.} В пакете реализованы три cost-функции,
описанные в~5.2:
\begin{itemize}
    \item \texttt{ClarkJacobStepneyNL(X, R)} — функция
        Кларка–Джейкоба–Степни с настраиваемыми параметрами
        $X, R$;
    \item \texttt{MinMaxWalsh()} — минимизация максимума спектра
        Уолша, эквивалентная $\mathrm{cost} = 256 - 2 \cdot N(S)$;
    \item \texttt{ThresholdNL(target)} — пороговая функция,
        штрафующая только нарушение порога нелинейности.
\end{itemize}
Все они приводятся к общему типу
\texttt{CostFunc = func([256]uint8) float64}, что позволяет
комбинировать алгоритмы и целевые функции независимо.

\textbf{Имитация отжига.} Функция \texttt{SA(state, sch)} принимает
состояние и структуру параметров \texttt{Schedule} с полями
$T_0, \alpha, M, T_{\min}, F_{\mathrm{target}}, \mathrm{MaxIter}$
и реализует алгоритм из подраздела~5.4. Принятые/отвергнутые
переходы учитываются в счётчиках, возвращаемых в составе
результата (\texttt{Accepted}, \texttt{AcceptedUp}) для последующей
диагностики. Опциональный callback \texttt{OnTick} вызывается в
начале каждого температурного уровня и позволяет журналировать
ход оптимизации без модификации основного цикла.

\textbf{Калибровка $T_0$.} Функция \texttt{EstimateT0(state, probes, p0)}
реализует алгоритм из подраздела~5.5: выполняет заданное число
пробных swap-переходов, накапливает среднее $\Delta F$ по
ухудшениям, после каждого пробного шага откатывается через
повторный swap. Возвращает значение $T_0$, согласованное с
целевой acceptance rate $p_0$.

\textbf{Параллельный multi-start.} Функция
\texttt{ParallelMultiStartSA(numRuns, sch, cf)} запускает
\texttt{numRuns} независимых SA в отдельных горутинах с помощью
\texttt{sync.WaitGroup}, синхронизированных только на завершении.
Каждая горутина строит собственное случайное начальное состояние
через \texttt{rand.Shuffle} и работает с собственным
\texttt{SAState}. Аналогичная функция \texttt{ParallelMultiStart}
реализована для hill climbing.

\textbf{Hill climbing.} Функция \texttt{HillClimb(state, maxIter)}
выполняет простейший локальный поиск: на каждой итерации
производится случайная транспозиция и проверяется условие
$F(S') < F(S)$; в противном случае выполняется откат через
повторный swap.

\subsection{Верификация инструмента модульным тестированием}

Корректность реализации обеспечивается набором автоматических
тестов, организованных в трёх слоях по возрастающей сложности
проверяемых объектов. Тесты используют встроенный в язык Go
механизм \texttt{testing} и выполняются командой \texttt{go test ./...}
без внешних зависимостей.

\textbf{Слой~1: модульные тесты базовых преобразований.}
В файлах \texttt{boolfunc/transforms\_test.go} и
\texttt{boolfunc/metrics\_test.go} проверяется корректность
реализованных алгоритмов на функциях малой размерности с известным
аналитическим ответом. В число эталонов входят константы $f \equiv 0$
и $f \equiv 1$, элементарные булевы функции $\mathrm{XOR}$ и
$\mathrm{AND}$ от двух переменных, линейные функции $f(x) = x_i$.
Для них вручную выписаны:
\begin{itemize}
    \item АНФ-коэффициенты (например, $f = x_1 \oplus x_2 \Rightarrow
        \mathrm{ANF} = [0, 1, 1, 0]$);
    \item значения спектра Уолша (например, для сбалансированной
        $f = x_1 \oplus x_2$: $W = [0, 0, 0, 4]$);
    \item значения нелинейности и алгебраической степени;
    \item значения алгебраической иммунности для линейной функции
        и констант.
\end{itemize}
Сравнение машинно вычисленных и аналитических значений выявляет
ошибки в реализации преобразований Мёбиуса, Уолша–Адамара и
гауссова исключения.

Дополнительно проверяются \emph{структурные} свойства преобразований:
инволютивность преобразования Мёбиуса, то есть $\mu(\mu(\mathrm{tv})) = \mathrm{tv}$;
обращение спектра Уолша в нуль на нулевой маске для сбалансированной
функции ($W_f(0) = 0 \Leftrightarrow w(f) = 2^{n-1}$); корректность
ленивого кэширования и изоляция возвращаемых копий от внутреннего
состояния структур.

\textbf{Слой~2: интеграционный тест на S-блоке AES.} Файл
\texttt{boolfunc/aes\_test.go} содержит таблицу подстановки AES
из стандарта~\cite{fips197} и три теста, использующие её как
комплексный эталон для модуля \texttt{boolfunc}:
\begin{itemize}
    \item \texttt{TestAES\_Balanced} — проверка, что все 255
        нетривиальных компонентных функций $f_\beta(x) = \langle\beta, S(x)\rangle$
        сбалансированы (вес $w(f_\beta) = 128$);
    \item \texttt{TestAES\_Nonlinearity} — проверка, что
        $\min_\beta N(f_\beta) = 112$;
    \item \texttt{TestAES\_AlgebraicDegree} — проверка, что
        $\max_\beta \deg(f_\beta) = 7$.
\end{itemize}
Прохождение этого набора означает корректность связки «вектор
истинности компонентной функции $\to$ Уолш/АНФ $\to$ метрика»
на содержательной задаче.

\textbf{Слой~3: тесты на эталонных S-блоках.} В файлах
\texttt{sbox/analysis\_test.go} и \texttt{sbox/ddt\_test.go}
проверяется совпадение всех четырёх метрик S-блока с табличными
значениями для четырёх стандартных подстановок, загружаемых из
JSON-файлов каталога \texttt{testdata/}. Эталонные значения
приведены в таблице~\ref{tab:reference-expected}.

\begin{table}[h]
\centering
\caption{Контрольные значения метрик эталонных S-блоков, проверяемые автоматическими тестами.}
\label{tab:reference-expected}
\begin{tabular}{lcccc}
\hline
S-блок & $N(S)$ & $\delta(S)$ & $\deg(S)$ & $AI(S)$ \\
\hline
AES (Rijndael)     & 112 & 4 & 7 & 3 \\
Camellia ($s_1$)   & 112 & 4 & 7 & 3 \\
SM4                & 112 & 4 & 7 & 3 \\
ГОСТ Р 34.12-2015 («Кузнечик») & 100 & 8 & 7 & 3 \\
\hline
\end{tabular}
\end{table}

Дополнительно файл \texttt{ddt\_test.go} содержит тест
\texttt{TestDDT\_properties}, выполняющий прямую машинную проверку
структурных свойств таблицы DDT (теорема о свойствах DDT биективного
S-блока из подраздела~3.2): пустота строки $\mathrm{DDT}[0]$,
сумма каждой строки $\sum_b \mathrm{DDT}[a][b] = 256$ для $a \neq 0$,
чётность всех элементов. Совпадение этих свойств для AES является
эмпирической верификацией теоретического утверждения.

Совокупное прохождение всех трёх слоёв тестов подтверждает
корректность реализованных алгоритмов: ошибки в преобразовании
Мёбиуса, FWHT, свёртке компонент, гауссовом исключении или
агрегации метрик S-блока приводят к расхождениям на одном из
эталонов с известным ответом.

\subsection*{ВЫВОДЫ}
\addcontentsline{toc}{subsection}{ВЫВОДЫ}

В шестой главе описан разработанный программный инструмент.
Архитектура состоит из четырёх независимых пакетов
(\texttt{boolfunc}, \texttt{sbox}, \texttt{generate},
\texttt{verify}) с однонаправленными зависимостями. Базовый
пакет \texttt{boolfunc} реализует булевы функции с поддержкой
ленивого вычисления АНФ, спектра Уолша и четырёх криптографических
метрик; пакет \texttt{sbox} строит над ним аппарат компонентных
функций, DDT и агрегированных метрик S-блока. Пакет
\texttt{generate} содержит реализацию имитации отжига с
калибровкой начальной температуры, hill climbing и параллельный
multi-start. Корректность модулей подтверждается сопоставлением
вычисленных метрик с табличными значениями для эталонных S-блоков
AES, Camellia, SM4 и «Кузнечик». Численные результаты приводятся
в следующей главе.


\newpage
%%%%%%%%%%%%%%%%%%%%%%%%%%%%%%%%%%%%%%%%%%%%%%%%%%%%%%%%%%%%%%%%%%%%%%%%%
% 7. Экспериментальные результаты
%%%%%%%%%%%%%%%%%%%%%%%%%%%%%%%%%%%%%%%%%%%%%%%%%%%%%%%%%%%%%%%%%%%%%%%%%
\section{Экспериментальные результаты}

% TODO: вступительный абзац — что исследуем в главе, какие гипотезы проверяем,
% какая аппаратная и программная среда. Источники чисел —
% knowledge-base/wiki/sa-experiments.md и sa-empirical-observations.md.

\subsection{Среда экспериментов и параметры алгоритмов}
% TODO: 16-ядерный вычислитель, версия Go, multi-start через горутины.
% Базовые параметры SA: $T_0$ из калибровки, $\alpha$, $M$, $T_{\min}$.

\subsection{Валидация на эталонных S-блоках}
% TODO: AES, Camellia, SM4, «Кузнечик». Сравнение вычисленных метрик
% с табличными значениями из главы~1. Контрольная точка корректности модуля.

\subsection{Бенчмарк случайных биективных подстановок}
% TODO: $N$ запусков Fisher-Yates. Распределение метрик: гистограммы или
% процентили для NL, $\delta$, $\deg$. Контрольная точка для SA.

\subsection{Одиночный SA: сравнение cost-функций}
% TODO: эксперименты E-04, E-05, E-06. MinMaxWalsh vs ThresholdNL vs CJS.
% Эмпирический потолок NL=100. Объяснение через структуру ландшафта.

\subsection{Multi-start и параллельные запуски}
% TODO: эксперимент E-08. Зависимость best-of-N от $N$. Формула $1-(1-p)^N$.

\subsection{Двухфазный алгоритм SA + HC}
% TODO: эксперименты E-07, E-09, E-10. Целевая глава. Достижение NL=112.

\subsection{Сравнительный анализ}
% TODO: финальная таблица random / single SA / SA+HC / эталоны (AES, Camellia, SM4, «Кузнечик»)
% по четырём метрикам (NL, $\delta$, $\deg$, AI).

\subsection*{ВЫВОДЫ}
\addcontentsline{toc}{subsection}{ВЫВОДЫ}
% TODO: 1–2 абзаца. Какие из заявленных в гл.~5 целевых значений достигнуты,
% какие — нет; что показал эксперимент в сравнении с теоретическими ожиданиями.


\newpage
%%%%%%%%%%%%%%%%%%%%%%%%%%%%%%%%%%%%%%%%%%%%%%%%%%%%%%%%%%%%%%%%%%%%%%%%%
% ЗАКЛЮЧЕНИЕ
%%%%%%%%%%%%%%%%%%%%%%%%%%%%%%%%%%%%%%%%%%%%%%%%%%%%%%%%%%%%%%%%%%%%%%%%%
\section*{ЗАКЛЮЧЕНИЕ}
\addcontentsline{toc}{section}{ЗАКЛЮЧЕНИЕ}

% TODO: 3–5 абзацев.
% 1. Краткое резюме поставленных задач (из ВВЕДЕНИЯ) и того, что сделано по каждой.
% 2. Главный численный результат: достигнутые значения метрик; сравнение с целевыми (AES: NL=112, DU=4, deg=7).
% 3. Где упёрлись и почему (стена NL=100 у single SA, пройденная через SA+HC).
% 4. Перспективы: parallel tempering, генетические алгоритмы, ключезависимые S-блоки.

\newpage
%%%%%%%%%%%%%%%%%%%%%%%%%%%%%%%%%%%%%%%%%%%%%%%%%%%%%%%%%%%%%%%%%%%%%%%%%
% СПИСОК СОКРАЩЁННЫХ ОБОЗНАЧЕНИЙ
%%%%%%%%%%%%%%%%%%%%%%%%%%%%%%%%%%%%%%%%%%%%%%%%%%%%%%%%%%%%%%%%%%%%%%%%%
\section*{СПИСОК СОКРАЩЁННЫХ ОБОЗНАЧЕНИЙ}
\addcontentsline{toc}{section}{СПИСОК СОКРАЩЁННЫХ ОБОЗНАЧЕНИЙ}
{
\setlength{\parindent}{0pt}
\begin{flushleft}
\begin{tabular}{@{}ll}
АНФ   & -- Алгебраическая нормальная форма \\
ПУА   & -- Преобразование Уолша--Адамара \\
ГОСТ  & -- Государственный стандарт \\
\mbox{} \\
\textit{AES}   & -- \textit{Advanced Encryption Standard}, стандарт шифрования \\
\textit{AI}    & -- Алгебраическая иммунность \textit{(Algebraic Immunity)} \\
\textit{APN}   & -- \textit{Almost Perfect Nonlinear}, почти совершенно нелинейная функция \\
\textit{DDT}   & -- \textit{Difference Distribution Table}, таблица разностного распределения \\
\textit{FWHT}  & -- \textit{Fast Walsh--Hadamard Transform}, быстрое преобразование Уолша--Адамара \\
\textit{GF}    & -- \textit{Galois Field}, конечное поле \\
\textit{HC}    & -- \textit{Hill Climbing}, восхождение к вершине \\
\textit{LAT}   & -- \textit{Linear Approximation Table}, таблица линейных аппроксимаций \\
\textit{MDS}   & -- \textit{Maximum Distance Separable}, код с максимально достижимым расстоянием \\
\textit{NL}    & -- Нелинейность \textit{(Nonlinearity)} \\
\textit{SA}    & -- \textit{Simulated Annealing}, имитация отжига \\
\textit{S-box} & -- \textit{Substitution box}, таблица замен \\
\textit{SPN}   & -- \textit{Substitution-Permutation Network}, подстановочно-перестановочная сеть \\
\textit{TV}    & -- \textit{Truth Vector}, вектор истинности \\
% TODO: дополнить по мере появления в тексте.
\end{tabular}
\end{flushleft}
}

\newpage
\section*{СПИСОК ИСПОЛЬЗОВАННЫХ ИСТОЧНИКОВ}
\addcontentsline{toc}{section}{СПИСОК ИСПОЛЬЗОВАННЫХ ИСТОЧНИКОВ}
\renewcommand{\refname}{\vspace*{-3em}}  % Уменьшаем пространство после заголовка

\begin{thebibliography}{9}
\renewcommand{\bibnumfmt}[1]{#1.}  % Убираем квадратные скобки из нумерации
\bibitem{biham1991}
Бихам, Э. Дифференциальный криптоанализ DES-подобных криптосистем [Текст] / Э. Бихам, А. Шамир. – Journal of Cryptology. – 1991. – Vol. 4, no. 1. – P. 3–72.

\bibitem{matsui1993}
Мацуи, М. Метод линейного криптоанализа для шифра DES [Текст] / М. Мацуи. – Advances in Cryptology – EUROCRYPT 1993. – LNCS 765. – Springer, 1994. – P. 386–397.

\bibitem{fips197}
NIST. FIPS 197: Advanced Encryption Standard (AES) [Текст]. – NIST, 2001. – 51 с.

\bibitem{gost2015}
ГОСТ Р 34.12-2015. Информационная технология. Криптографическая защита информации. Блочные шифры [Текст]. – М.: Стандартинформ, 2015.

\bibitem{biryukov2016}
Бирюков, А. Реверс-инжиниринг S-блока Стрибога, Кузнечика и STRIBOBr1 [Текст] / А. Бирюков, Л. Перрин, А. Удовенко. – Advances in Cryptology – EUROCRYPT 2016. – LNCS 9665. – Springer, 2016. – P. 372–402.

\bibitem{serpent1998}
Андерсон, Р. Serpent: предложение для стандарта AES [Текст] / Р. Андерсон, Э. Бихам, Л. Кнудсен. – AES Submission Document. – 1998. – 26 с.

\bibitem{twofish1998}
Шнайер, Б. Twofish: 128-битный блочный шифр [Текст] / Б. Шнайер и др. – AES Submission Document. – 1998. – 137 с.

\bibitem{sm4_2006}
OSCCA. GM/T 0002-2012. SM4 Block Cipher Algorithm [Текст]. – 2012.

\bibitem{camellia2001}
Аоки, К. Camellia: 128-битный блочный шифр, подходящий для разных платформ. Проектирование и анализ [Текст] / К. Аоки и др. – Selected Areas in Cryptography. – LNCS 2012. – Springer, 2001. – P. 39–56.

\bibitem{lightweight_cryptographic}
АНСАРИ, Ш. А., АЛИ, С. A systematic review of lightweight cryptographic schemes for security and privacy in IoT [Текст] // Discover Computing. – 2025. – Vol. 28, no. 1. – P. 1–57. – DOI: 10.1007/s10791-025-09755-3.

\bibitem{alferov2002}
Алферов, А. П. Основы криптографии: учебное пособие [Текст] / А. П. Алферов, А. Ю. Зубов, А. С. Кузьмин, А. В. Черемушкин. – 2-е изд., испр. и доп. – М.: Гелиос АРВ, 2002. – 480 с.

\bibitem{nyberg1993apn}
Nyberg, K. Differentially uniform mappings for cryptography [Text] / K. Nyberg // Advances in Cryptology – EUROCRYPT 1993. – LNCS 765. – Springer, 1994. – P. 55–64.

\bibitem{carlet2021book}
Carlet, C. Boolean Functions for Cryptography and Coding Theory [Text] / C. Carlet. – Cambridge: Cambridge University Press, 2021. – 562 p.

\bibitem{lai1994higher}
Lai, X. Higher Order Derivatives and Differential Cryptanalysis [Text] / X. Lai // Communications and Cryptography: Two Sides of One Tapestry. – Boston: Springer, 1994. – P. 227–233.

\bibitem{knudsen1995truncated}
Knudsen, L. R. Truncated and Higher Order Differentials [Text] / L. R. Knudsen // Fast Software Encryption (FSE 1995). – LNCS 1008. – Springer, 1995. – P. 196–211.

\bibitem{courtois2003algebraic}
Courtois, N. Algebraic Attacks on Stream Ciphers with Linear Feedback [Text] / N. Courtois, W. Meier // Advances in Cryptology – EUROCRYPT 2003. – LNCS 2656. – Springer, 2003. – P. 345–359.

\bibitem{meier2004algebraic}
Meier, W. Algebraic Attacks and Decomposition of Boolean Functions [Text] / W. Meier, E. Pasalic, C. Carlet // Advances in Cryptology – EUROCRYPT 2004. – LNCS 3027. – Springer, 2004. – P. 474–491.

\bibitem{kirkpatrick1983sa}
Kirkpatrick, S. Optimization by Simulated Annealing [Text] / S. Kirkpatrick, C. D. Gelatt, M. P. Vecchi // Science. – 1983. – Vol. 220, no. 4598. – P. 671–680.

\bibitem{cerny1985sa}
Černý, V. Thermodynamical Approach to the Traveling Salesman Problem: An Efficient Simulation Algorithm [Text] / V. Černý // Journal of Optimization Theory and Applications. – 1985. – Vol. 45, no. 1. – P. 41–51.

\bibitem{gemangeman1984sa}
Geman, S. Stochastic Relaxation, Gibbs Distributions, and the Bayesian Restoration of Images [Text] / S. Geman, D. Geman // IEEE Transactions on Pattern Analysis and Machine Intelligence. – 1984. – Vol. PAMI-6, no. 6. – P. 721–741.

\bibitem{benameur2004t0}
Ben-Ameur, W. Computing the Initial Temperature of Simulated Annealing [Text] / W. Ben-Ameur // Computational Optimization and Applications. – 2004. – Vol. 29, no. 3. – P. 369–385.

\bibitem{marti2013multistart}
Martí, R. Multi-Start Methods for Combinatorial Optimization [Text] / R. Martí, M. G. C. Resende, C. C. Ribeiro // European Journal of Operational Research. – 2013. – Vol. 226, no. 1. – P. 1–8.

\bibitem{clark2005sboxes}
Clark, J. A. The Design of S-boxes by Simulated Annealing [Text] / J. A. Clark, J. L. Jacob, S. Stepney // New Generation Computing. – 2005. – Vol. 23, no. 3. – P. 219–231.

\bibitem{mdpi2022sboxes}
Optimization of a Simulated Annealing Algorithm for S-Boxes Generating [Text] // Sensors (MDPI). – 2022. – Vol. 22, no. 16. – P. 6073. – DOI: 10.3390/s22166073.

\end{thebibliography}


\newpage
%%%%%%%%%%%%%%%%%%%%%%%%%%%%%%%%%%%%%%%%%%%%%%%%%%%%%%%%%%%%%%%%%%%%%%%%%
% ПРИЛОЖЕНИЕ А
%%%%%%%%%%%%%%%%%%%%%%%%%%%%%%%%%%%%%%%%%%%%%%%%%%%%%%%%%%%%%%%%%%%%%%%%%
\section*{ПРИЛОЖЕНИЕ А}
\begin{center}
    \textbf{Программная реализация инструмента анализа и генерации S-блоков}
\end{center}
\addcontentsline{toc}{section}{ПРИЛОЖЕНИЕ А \phantom{M}Программная реализация инструмента анализа и генерации S-блоков}

% TODO: подключить листинги ключевых Go-файлов через \codefromfile{file}{language}.
% Файлы предварительно копируются в Listings/ из соответствующих пакетов проекта.
% Примеры (раскомментировать по мере готовности):

% \subsection*{boolfunc/metrics.go}
% \codefromfile{metrics.go}{go}
%
% \subsection*{sbox/ddt.go}
% \codefromfile{ddt.go}{go}
%
% \subsection*{generate/annealing.go}
% \codefromfile{annealing.go}{go}
%
% \subsection*{generate/cost.go}
% \codefromfile{cost.go}{go}
